% =====================================================================
%  Lecture 2026-10-15
%
%  Normalized silently:
%   - H_1, H_2 and H for Hilbert spaces, as in Section 14
%   - the supremum variable in the norm identity of the adjoint, called
%     z in the source as well as the representing vector, renamed u
%   - in property (5), <Tx, Tx> = <x, T*Tx> directly from the definition
%   - "real space" and "complex space" for L_{2,R} and L_{2,C}, as in
%     Example 7.8
%   - the proof of the product theorem placed directly after its
%     statement, before the remark it precedes in the source
%
%  CORRECTED SOURCE ERRORS
%   [p1] lemma hypothesis T in C(H_1, H_2)  ->  B(H_1, H_2)
%   [p2] "if T, S are adjoint"  ->  self-adjoint
%   [p2] product theorem proof (TS)* = S*T* = ST = TS, assuming the
%        claim  ->  (TS)* = ST, so (TS)* = TS iff ST = TS
%   [p2] ||U|| = 1 for every Hilbert space  ->  H != {0}
%   [p3] Fourier systems without index sets, as families  ->
%        enumerated as sequences e_1, e_2, ..., in the order of
%        Example 7.8; the complex system over all k in Z
%   [p3] index clash: the real system starts at e_0, the ONB definition
%        at e_1  ->  every orthonormal sequence indexed from 1
%
%  ADDITIONS
%   [p1] remark ||v|| = sup_{||u|| <= 1} |<v, u>|, with proof; the
%        source's "properties of the inner product"
%   [p1] adjoint theorem: f linear; uniqueness in Theorem 14.1 cited for
%        well-definedness; T* linear; the uniqueness of the adjoint
%        stated as a remark
%   [p1] conjugate transpose example: the index computation
%   [p1] properties (1) to (3): one-line proofs, where the source says
%        trivial; (4) and (5): Proposition 9.8 cited, the step to
%        equality of operators via the uniqueness remark, and
%        ||TT*|| = ||T||^2 derived from the first part applied to T*
%   [p1] unitary: T invertible, so that T^{-1} is bounded, as in
%        Definition 10.2
%   [p1] lemma (1): the identification of the supremum with ||T|| cited;
%        lemma (2): the trivial direction, and the step from the two
%        equations to <Tx, y> = 0
%   [p1] rotation counterexample: the formula of the rotation
%   [p2] <x, Tx> = <T*x, x>, justified; Lemma 16.2 cited in (2)
%   [p2] unitary (5): the direction "unitary => surjective and
%        isometric", from (1); T^{-1} = T* bounded
%   [p2] unitary group remark: identity unitary, closure from (3), (4)
%   [p2] shift example: S isometric, e_1 not in im S, S*S = I and
%        SS* e_1 = 0, written out
%   [p3] orthonormality of the trigonometric systems, by integration
%   [p3] the ONB definition requires (e_n) orthonormal
%
%  REPORTED BUT UNMODIFIED GAPS
%   [p2] in finite dimensions, unitary iff isometric, asserted
%   [p3] whether the trigonometric systems are orthonormal bases is not
%        addressed here
%
%  OMITTED
%   [p1] course logistics in the top right margin
%   [p2] a doodle
% =====================================================================

% [p1]
\section{Adjoint operators}

\begin{remark}\label{rem:norm-sup}
In an inner product space, $\norm{v} = \sup_{\norm{u} \le 1} \abs{\ip{v}{u}}$
for every $v$. Indeed $\abs{\ip{v}{u}} \le \norm{v}$ for $\norm{u} \le 1$ by
the Cauchy--Schwarz inequality, Theorem~\ref{thm:cauchy-schwarz}, and for
$v \ne 0$ the unit vector $u = v / \norm{v}$ gives $\ip{v}{u} = \norm{v}$.
\end{remark}

\begin{definition}[Adjoint operator]\label{def:adjoint}
Let $T \in \Bop{H_1}{H_2}$. The \emph{adjoint} of $T$ is the map
$T^* \colon H_2 \to H_1$ given by
\[
  \ip{Tx}{y} = \ip{x}{T^* y} \qquad \text{for all } x \in H_1,\ y \in H_2 .
\]
\end{definition}

\begin{theorem}\label{thm:adjoint}
$T^*$ is well defined, $T^* \in \Bop{H_2}{H_1}$, and $\norm{T} = \norm{T^*}$.
\end{theorem}

\begin{proof}
Fix $y \in H_2$ and define $f(x) = \ip{Tx}{y}$ for $x \in H_1$. It is linear,
since $T$ is linear and the inner product is linear in its first variable.
By the Cauchy--Schwarz inequality,
\[
  \sup_{\norm{x} \le 1} \abs{f(x)} = \sup_{\norm{x} \le 1} \abs{\ip{Tx}{y}}
  \le \sup_{\norm{x} \le 1} \norm{Tx}\,\norm{y}
  \le \norm{T}\,\norm{y} < \infty ,
\]
so $f \in H_1'$. By the Riesz representation, Theorem~\ref{thm:riesz-hilbert},
there is a unique $z \in H_1$ with $f(x) = \ip{x}{z}$ for all $x \in H_1$.
Then $\ip{Tx}{y} = \ip{x}{z}$, and by the uniqueness of $z$ the definition
$T^* y = z$ is meaningful.

$T^*$ is linear: for $y, y' \in H_2$ and $\beta \in \K$, by
Remark~\ref{rem:conj-linear},
\[
  \ip{x}{T^*(y + \beta y')} = \ip{Tx}{y + \beta y'}
  = \ip{Tx}{y} + \bar{\beta}\ip{Tx}{y'}
  = \ip{x}{T^* y + \beta T^* y'}
\]
for all $x \in H_1$, and the uniqueness in Theorem~\ref{thm:riesz-hilbert}
gives $T^*(y + \beta y') = T^* y + \beta T^* y'$.

Finally, by Remark~\ref{rem:norm-sup},
\[
  \norm{T} = \sup_{\norm{x} \le 1} \norm{Tx}
  = \sup_{\substack{\norm{x} \le 1\\ \norm{u} \le 1}} \abs{\ip{Tx}{u}}
  = \sup_{\substack{\norm{x} \le 1\\ \norm{u} \le 1}} \abs{\ip{x}{T^* u}}
  = \sup_{\norm{u} \le 1} \norm{T^* u} = \norm{T^*} ,
\]
where the second and the fourth equalities are Remark~\ref{rem:norm-sup},
applied to $Tx$ and to $T^* u$. Therefore $T^*$ is also bounded.
\end{proof}

\begin{remark}\label{rem:adjoint-unique}
If $R \colon H_2 \to H_1$ satisfies $\ip{Tx}{y} = \ip{x}{Ry}$ for all
$x \in H_1$ and $y \in H_2$, then $R = T^*$: for fixed $y$,
$\ip{x}{Ry - T^* y} = 0$ for all $x$, and $x = Ry - T^* y$ gives
$Ry = T^* y$.
\end{remark}

\begin{example}\label{ex:conj-transpose}
Let $T \in \operatorname{End}(\C^n)$, identified with its matrix $(t_{ij})$ in
the standard basis, and let $\C^n$ carry the inner product
$\ip{x}{y} = \sum_{i=1}^{n} x_i \overline{y_i}$. Then
$T^* \in \operatorname{End}(\C^n)$, and $T^* = \overline{T}^{\mathsf{T}}$,
the conjugate transpose, since
\[
  \ip{Tx}{y} = \sum_{i=1}^{n} \sum_{j=1}^{n} t_{ij} x_j \overline{y_i}
  = \sum_{j=1}^{n} x_j\, \overline{\sum_{i=1}^{n} \overline{t_{ij}}\, y_i}
  = \ip[\big]{x}{\overline{T}^{\mathsf{T}} y} .
\]
\end{example}

\begin{theorem}[Properties of the adjoint]\label{thm:adjoint-properties}
Let $S, T \in \Bop{H_1}{H_2}$ and $\alpha \in \K$. Then:
\begin{enumerate}
  \item $(S + T)^* = S^* + T^*$;
  \item $(\alpha T)^* = \bar{\alpha} T^*$;
  \item $(T^*)^* = T$;
  \item if $H_1 = H_2$, then $(ST)^* = T^* S^*$;
  \item $\norm{T^* T} = \norm{T T^*} = \norm{T}^2$.
\end{enumerate}
\end{theorem}

\begin{proof}
Items (1) to (3) follow from Remark~\ref{rem:adjoint-unique}, since the
right-hand side satisfies the defining identity of the left-hand side:
\begin{enumerate}
  \item $\ip{x}{(S^* + T^*)y} = \ip{Sx}{y} + \ip{Tx}{y} = \ip{(S + T)x}{y}$;
  \item $\ip{x}{\bar{\alpha} T^* y} = \alpha\ip{x}{T^* y}
    = \alpha\ip{Tx}{y} = \ip{(\alpha T)x}{y}$;
  \item $\ip{T^* y}{x} = \overline{\ip{x}{T^* y}} = \overline{\ip{Tx}{y}}
    = \ip{y}{Tx}$ for $y \in H_2$, $x \in H_1$, which is the defining
    identity of $(T^*)^*$ with $T$ in its place.
\end{enumerate}
(4) For all $x, y$,
$\ip{x}{(ST)^* y} = \ip{STx}{y} = \ip{Tx}{S^* y} = \ip{x}{T^* S^* y}$, so
$(ST)^* = T^* S^*$ by Remark~\ref{rem:adjoint-unique}.

(5) By the Cauchy--Schwarz inequality,
\[
  \begin{aligned}
    \norm{T}^2 &= \sup_{\norm{x} \le 1} \norm{Tx}^2
    = \sup_{\norm{x} \le 1} \ip{Tx}{Tx}
    = \sup_{\norm{x} \le 1} \ip{x}{T^* T x} \\
    &\le \sup_{\norm{x} \le 1} \norm{T^* T x}\,\norm{x}
    \le \sup_{\norm{x} \le 1} \norm{T^* T x} = \norm{T^* T} ,
  \end{aligned}
\]
and by Proposition~\ref{prop:composition} and Theorem~\ref{thm:adjoint},
$\norm{T^* T} \le \norm{T^*}\,\norm{T} = \norm{T}^2$. So
$\norm{T^* T} = \norm{T}^2$. Applied to $T^*$ in place of $T$, with (3) and
$\norm{T^*} = \norm{T}$, this gives
$\norm{T T^*} = \norm{(T^*)^* T^*} = \norm{T^*}^2 = \norm{T}^2$.
\end{proof}

\section{Self-adjoint, unitary and normal operators}

\begin{definition}\label{def:operator-classes}
An operator $T \in \Bopone{H}$ is called
\begin{enumerate}
  \item \emph{self-adjoint}, if $T = T^*$;
  \item \emph{unitary}, if $T$ is invertible and $T^* = T^{-1}$;
  \item \emph{normal}, if $T^* T = T T^*$.
\end{enumerate}
\end{definition}

\begin{lemma}\label{lem:zero-operator}
Let $T \in \Bop{H_1}{H_2}$. Then:
\begin{enumerate}
  \item $T = 0$ if and only if $\ip{Tx}{y} = 0$ for all $x \in H_1$ and
    $y \in H_2$;
  \item if $H_1 = H_2 = H$ is a complex Hilbert space, then $T = 0$ if and
    only if $\ip{Tx}{x} = 0$ for all $x \in H$.
\end{enumerate}
\end{lemma}

\begin{remark}
Item (2) does not hold on real spaces. For example, let $T$ be the rotation
of $\R^2$ by $90^\circ$, $T(x_1, x_2) = (-x_2, x_1)$. Then $T \ne 0$, but
$\ip{Tx}{x} = -x_2 x_1 + x_1 x_2 = 0$ for every $x$.
\end{remark}

\begin{proof}[Proof of Lemma~\ref{lem:zero-operator}]
(1) If $T = 0$, then $\ip{Tx}{y} = \ip{0}{y} = 0$. Conversely, if
$\ip{Tx}{y} = 0$ for all $x, y$, then
$\sup_{\norm{x} \le 1, \norm{y} \le 1} \abs{\ip{Tx}{y}} = 0$. By the first two
equalities in the proof of Theorem~\ref{thm:adjoint} this supremum is
$\norm{T}$, so $\norm{T} = 0$ and $T = 0$.

(2) If $T = 0$, then $\ip{Tx}{x} = 0$. Conversely, assume $\ip{Tx}{x} = 0$
for all $x \in H$, with $H$ complex. Let $x, y \in H$, $\alpha \in \C$, and
write $v = \alpha x + y$. Then
\[
  0 = \ip{Tv}{v} = \ip{T(\alpha x + y)}{\alpha x + y}
  = \abs{\alpha}^2 \underbrace{\ip{Tx}{x}}_{0}
    + \underbrace{\ip{Ty}{y}}_{0}
    + \alpha\ip{Tx}{y} + \bar{\alpha}\ip{Ty}{x} .
\]
% [p2]
Now $\alpha = 1$ gives $\ip{Tx}{y} + \ip{Ty}{x} = 0$, and $\alpha = i$ gives
$i\ip{Tx}{y} - i\ip{Ty}{x} = 0$. Dividing the second equation by $i$ and
adding it to the first, $2\ip{Tx}{y} = 0$. So $\ip{Tx}{y} = 0$ for all
$x, y$, and $T = 0$ by (1).
\end{proof}

\begin{theorem}\label{thm:self-adjoint-real}
Let $T \in \Bopone{H}$. Then:
\begin{enumerate}
  \item if $T = T^*$, then $\ip{Tx}{x} \in \R$ for all $x \in H$;
  \item if $H$ is complex, then $T = T^*$ if and only if
    $\ip{Tx}{x} \in \R$ for all $x \in H$.
\end{enumerate}
\end{theorem}

\begin{proof}
For every $T \in \Bopone{H}$ and $x \in H$,
$\ip{T^* x}{x} = \overline{\ip{x}{T^* x}} = \overline{\ip{Tx}{x}}
= \ip{x}{Tx}$.

(1) If $T = T^*$, then
$\overline{\ip{Tx}{x}} = \ip{x}{Tx} = \ip{T^* x}{x} = \ip{Tx}{x}$, so
$\ip{Tx}{x} \in \R$.

(2) One direction is (1). Assume $\ip{Tx}{x} \in \R$ for all $x$. Then
$\ip{Tx}{x} = \overline{\ip{Tx}{x}} = \ip{x}{Tx} = \ip{T^* x}{x}$, so
\[
  0 = \ip{Tx}{x} - \ip{T^* x}{x} = \ip{(T - T^*)x}{x}
  \qquad \text{for all } x \in H .
\]
By Lemma~\ref{lem:zero-operator} (2), $T - T^* = 0$, that is, $T = T^*$.
\end{proof}

\begin{theorem}\label{thm:product-self-adjoint}
Let $S, T \in \Bopone{H}$ be self-adjoint. Then $TS$ is self-adjoint, that
is $(TS)^* = TS$, if and only if $TS = ST$.
\end{theorem}

\begin{proof}
By Theorem~\ref{thm:adjoint-properties} (4), $(TS)^* = S^* T^* = ST$. So
$(TS)^* = TS$ if and only if $ST = TS$.
\end{proof}

\begin{remark}
If $T$ and $S$ are self-adjoint, then so is $T + S$, as
$(T + S)^* = T^* + S^* = T + S$. Also, for $\alpha \in \R$,
$(\alpha T)^* = \bar{\alpha} T^* = \alpha T$.
\end{remark}

\begin{theorem}[Unitary operators]\label{thm:unitary}
Let $H \ne \set{0}$, and let $U, V \in \Bopone{H}$ be unitary. Then:
\begin{enumerate}
  \item $U$ is \emph{isometric}, that is, $\norm{Ux} = \norm{x}$ for all
    $x \in H$;
  \item $\norm{U} = 1$;
  \item $U^{-1} = U^*$ is unitary;
  \item $UV$ is also unitary;
  \item on a complex Hilbert space $H$, an operator $T \in \Bopone{H}$ is
    unitary if and only if $T$ is surjective and isometric.
\end{enumerate}
\end{theorem}

\begin{proof}
(1) As $U^* U = U^{-1} U = I$,
$\norm{Ux}^2 = \ip{Ux}{Ux} = \ip{x}{U^* U x} = \ip{x}{x} = \norm{x}^2$.

(2) By (1), $\norm{U} = \sup_{\norm{x} \le 1} \norm{Ux}
= \sup_{\norm{x} \le 1} \norm{x} = 1$, the last supremum being $1$ because
$H \ne \set{0}$.

(3) $U^*$ is invertible, with bounded inverse $U$. It is unitary if
$(U^*)^* = (U^*)^{-1}$, that is, by
Theorem~\ref{thm:adjoint-properties} (3), if $U = (U^*)^{-1}$, which holds
because $U^{-1} = U^*$.

(4) $UV$ is invertible, with bounded inverse $V^{-1} U^{-1}$, and
$(UV)^* = V^* U^* = V^{-1} U^{-1} = (UV)^{-1}$.

(5) If $T$ is unitary, it is isometric by (1), and surjective, being
invertible. Conversely, let $T$ be surjective and isometric. It is injective,
since $Tx = 0$ gives $\norm{x} = \norm{Tx} = 0$, so $x = 0$; with the
assumption, $T$ is bijective. For every $x$,
\[
  \ip{T^* T x}{x} = \overline{\ip{x}{T^* T x}} = \ip{Tx}{Tx}
  = \norm{Tx}^2 = \norm{x}^2 = \ip{x}{x} ,
\]
so $\ip{(T^* T - I)x}{x} = 0$, and $T^* T = I$ by
Lemma~\ref{lem:zero-operator} (2), since the space is complex. Moreover
$T T^* = T T^* T T^{-1} = T T^{-1} = I$. So $T^*$ is a two-sided inverse of
$T$, $T^{-1} = T^*$ is bounded, and $T$ is unitary.
\end{proof}

\begin{remark}
The unitary operators form a group under composition, the \emph{unitary
group}: $I$ is unitary, and by (3) and (4) inverses and products of unitary
operators are unitary.
\end{remark}

\begin{remark}
In finite dimensions, $T$ is unitary if and only if $T$ is isometric.
\end{remark}

\begin{example}[Shift operator]\label{ex:shift}
The surjectivity assumption in Theorem~\ref{thm:unitary} (5) is necessary
in infinite dimensions. Take the shift operator
\[
  S \colon \lp{2} \to \lp{2}, \qquad
  (x_1, x_2, \dots) \longmapsto (0, x_1, x_2, \dots) .
\]
$S$ is isometric, since $\norm{Sx}^2 = \sum_{i} \abs{x_i}^2 = \norm{x}^2$,
but it is not surjective, since every element of $\operatorname{im} S$ has
first coordinate $0$, so $e_1 \notin \operatorname{im} S$. $S$ is not
unitary either. Its adjoint is given by the “chop”
\[
  S^* \colon (y_1, y_2, \dots) \longmapsto (y_2, y_3, \dots) .
\]
This really is the adjoint of $S$:
\[
  \begin{aligned}
    \ip{Sx}{y} &= \ip{(0, x_1, x_2, \dots)}{(y_1, y_2, \dots)}
    = x_1 \overline{y_2} + x_2 \overline{y_3} + \dots \\
    &= \ip{(x_1, x_2, \dots)}{(y_2, y_3, \dots)} = \ip{x}{S^* y} .
  \end{aligned}
\]
Then $S^* S = I$, but $S S^* \ne I$, as $S S^* e_1 = S 0 = 0$. So $S$ is not
unitary.
\end{example}

% [p3]
\section{Orthonormal sequences}

\begin{definition}[Orthonormal sequence]\label{def:orthonormal}
A sequence $(e_n)$ in an inner product space is \emph{orthonormal} if
$\ip{e_i}{e_j} = 0$ for all $i \ne j$, and $\ip{e_i}{e_i} = 1$ for all $i$.
\end{definition}

\begin{example}\label{ex:trig-orthonormal}
The Fourier systems of Example~\ref{ex:fourier}, normalized, are orthonormal
sequences. In the real space $\Lp{2}[0,2\pi]$, take
\[
  e_1 = \frac{1}{\sqrt{2\pi}} , \qquad
  e_{2k}(t) = \frac{1}{\sqrt{\pi}} \cos(kt) , \qquad
  e_{2k+1}(t) = \frac{1}{\sqrt{\pi}} \sin(kt) , \qquad k \ge 1 .
\]
In the complex space $\Lp{2}[0,2\pi]$, take $\frac{1}{\sqrt{2\pi}} e^{ikt}$,
$k \in \Z$, in the order $k = 0, 1, -1, 2, -2, \dots$:
\[
  e_1 = \frac{1}{\sqrt{2\pi}} , \qquad
  e_{2k}(t) = \frac{1}{\sqrt{2\pi}} e^{ikt} , \qquad
  e_{2k+1}(t) = \frac{1}{\sqrt{2\pi}} e^{-ikt} , \qquad k \ge 1 .
\]
Orthonormality follows from $\int_0^{2\pi} e^{i(k-l)t} \dd t = 2\pi$ for
$k = l$ and $0$ for $k \ne l$; in the real case, from the product-to-sum
formulas, with $\int_0^{2\pi} \cos^2(kt) \dd t
= \int_0^{2\pi} \sin^2(kt) \dd t = \pi$ for $k \ge 1$.
\end{example}

\begin{definition}[Orthonormal basis]\label{def:onb}
An orthonormal sequence $(e_n)$ in a Hilbert space $H$ is an
\emph{orthonormal basis} (ONB) if $\operatorname{span}\set{e_1, e_2, \dots}$
is dense in $H$.
\end{definition}
